\section{Challenges, Risks, and Limitations}

\subsection{Shared Project Risks}

\begin{itemize}
  \item Real school library data may not be available before the midterm, so some validation must use simulated data or sample images.
  \item Module integration depends on consistent field names and data formats across frontend, backend, database, and image algorithm work.
  \item The final system must avoid committing passwords, API keys, personal data, or confidential images to the repository.
  \item A live demo is not required for the midterm, but the submitted evidence must still make implemented work verifiable.
\end{itemize}

\subsection{Data Structure Risks}

% Data structure teammate: add database-specific risks, such as schema changes,
% import failures, data volume, indexing, or mismatch between simulated and real data.

\subsection{Frontend Risks}

% Frontend teammate: add UI/backend integration risks, browser compatibility,
% incomplete workflows, API availability, or user experience limitations.

\subsection{Image Algorithm Risks}

The image algorithm module has three main limitations. First, the team does not yet have official school library video data with ground-truth labels. Therefore, the current module can be verified as a working pipeline, but not yet as a formally accurate model. Metrics such as precision, recall, and false-positive rate should be reported only after labeled real frames are collected.

Second, ROI calibration is critical. If the ROI is too small, a real person or object may fall outside the seat area and be missed. If the ROI is too large, objects from neighboring seats may be assigned incorrectly. The long-term design should use one calibrated ROI per physical seat for each camera view, with optional refinements if a layout is especially complex.

Third, YOLOv8n is a general-purpose COCO pre-trained model. It is lightweight and suitable for a CPU midterm baseline, but it is not specialized for library seat occupation. Final improvement may require threshold tuning, a larger detector, or fine-tuning with labeled library images.
